% Source: https://artofproblemsolving.com/community/c3h1449780_what_is_the_correct_answer_
% Stable identity: aops-topic-1449780-q01
Let $f(x)=x^6-3x^5-6x^3-x+8$.

\textbf{No negative roots.} If $x<0$, then $x^6>0$, $-3x^5>0$, $-6x^3>0$, $-x>0$ and $8>0$, so $f(x)>0$. Also $f(0)=8\ne0$. This rules out (B) and (C).

\textbf{Positive roots.} $f(1)=1-3-6-1+8=-1<0$ while $f(0)=8>0$, so by the intermediate value theorem there is a root in $(0,1)$. (Since $f(x)\to+\infty$ as $x\to\infty$, there is another root greater than $1$; by Descartes' rule of signs there are exactly two positive roots.) This rules out (A).

The answer is $\boxed{\textbf{(D)}}$.
